% Modernised reading — fonds Alexandre Grothendieck, shelfmark 63, pages 1-100.
%
% This is an INTERPRETATION, not a transcription. It re-reads the pages in
% today's notation and vocabulary, states the hypotheses the manuscript leaves
% implicit, and drops the critical apparatus so that the mathematics reads
% straight through. Everything the brackets and underlines carried moves into a
% footnote. It is held to being mathematically correct as it stands: where the
% page is loose or mistaken, the true statement is what appears here and the
% footnote says what the page has. In any dispute of substance the
% transcription governs: whoever cites Grothendieck must cite batch-01.fr.tex
% to batch-05.fr.tex.
%
% Pass: Opus 5.5 (claude-opus-5-5), 2026-10-03 — first pass, unchecked against
% the pages by a human. Not measured against the Opus 5 reference reading of
% folder 115. The folder was read whole — five batches, pages 1-100 — before
% any of this was written, and it is written as ONE file for the shelfmark.
%
% What the folder is. A « formulaire »: most of its hundred pages are other
% people's texts, which the transcriptions do not re-set and neither does this
% reading — J. Tate's letter to Cassels (pp. 2-13), P. Deligne's typescript
% « Courbes elliptiques : formulaire d'après J. Tate » (pp. 15-33), five older
% leaves of Tate (pp. 35-39), two typescripts of A. O. L. Atkin (pp. 40-63), a
% typed list of the CM curves with rational j (pp. 79-80), an anonymous
% bibliography and a « traduction libre d'un manuscrit inédit de Tate »
% (pp. 82-93). They are described, not restated. His own mathematics is in
% three places: the pencilled questions on Deligne's typescript (pp. 16-33,
% attribution uncertain, as the transcriptions say), seven leaves on the
% periods and the two splittings of H^1_dR (pp. 65-76), and five pages on the
% modular stack near j = 0, 12^3 and on level 3 (pp. 95-99, two of them a
% heavily struck palimpsest).
%
% NOTATION fixed for the whole folder (stated in the spine section):
%   E an elliptic curve, L = \underline{\omega} its line of invariant
%     differentials (his underlined omega), omega a basis of L;
%   V = H^1_dR(E/C), with 0 -> L -> V -> L^{-1} -> 0;
%   eta = x dx/2y (his eta), alpha, beta the scalars of p. 66
%     (eta + alpha omega = beta omega-bar);
%   A in L-bar^{-1} (x) L and B in (L-bar L)^{-1}: the maps of pp. 72-76,
%     which he ALSO calls alpha and beta — renamed here, footnoted;
%   varpi = alpha omega^{(x)2} the canonical section of p. 66 (his pi, renamed
%     because pi is Deligne's invariant differential and the number pi);
%   Lambda = Z omega_1 + Z omega_2, tau = omega_2/omega_1, Im tau > 0;
%   G_2(tau) = sum_n sum'_m (m + n tau)^{-2}, G_2^*(tau) = G_2(tau) - pi/Im tau;
%   M the « site modulaire » of pp. 31 and 95-99, read as the stack of
%     Weierstrass curves without additive fibres (compactified M_{1,1}), so
%     that M^infty = div(Delta) makes sense; M' = V(c_4), M'' = V(c_6);
%   M_3 the level-3 moduli (full level 3 structure), G = GL(2, F_3),
%     G' = SL(2, F_3), B, B' the Borel subgroups;
%   j is the USUAL invariant c_4^3/Delta throughout (j = 12^3 for the curve
%     with four automorphisms) — not the j = 1 - J of folder 62, pp. 18-53.
%
% Substantive departures from the pages, each footnoted where it occurs:
%   - p. 65: the series for y runs over a != 0 (the page's sum has no prime);
%   - p. 68: the sum in eta_2 excludes a = omega_2 (page: a != omega_1);
%   - p. 70: the factor 3 of -3(|omega_1|^2 - |omega_2|^2)/(omega_1 omega_2)
%     is lost from the third line on; the closed form alpha =
%     omega_1^{-2} G_2^*(tau), obtained here from Legendre's relation, is ours
%     and replaces the unfinished series;
%   - p. 74/76: « 4 rho sigma = 1 - alpha alpha-bar » has the wrong sign: the
%     identity is 4 rho sigma = A A-bar - 1 (with eta normalised by
%     int omega ^ eta = 1); sigma has no sign, and « 0 <= alpha alpha-bar <= 1 »
%     is false (A = 0 at tau = i, |A| -> infinity as Im tau -> infinity);
%   - p. 99: the 2-Sylow of SL(2, F_3) is of ORDER 8 (index 3), not « d'indice
%     8 »;
%   - p. 79 (pencil): « les d sont premiers » holds for d > 1, not for d = 1.
%
% Announced and never established: the expression of (arg A)^6 « en fonction
% de j » (p. 76) — answered here in a footnote, not by the page; the answers to
% the questions of p. 66 and p. 33; the end of the sentence « Donc les
% quatre ... » (p. 97), which p. 99 resumes; the structure of the complete
% local ring at the supersingular points and the equations of c_4, c_6 there
% (p. 98, « Question »); which of the two surjections G' -> F_4^* occurs
% (p. 99, « Il faudrait encore déterminer lequel c'est »).
\documentclass[11pt,a4paper]{article}
\input{../preamble/grothendieck}

\folder{63}
\pages{1}{100}
\dating{[à partir de 1964-vers 1970]}
\watermark{Édition de démonstration\\interprétation personnelle de l'œuvre}
\foldertitle{Formulaire courbes elliptiques : notes et copies de notes manuscrites (s.d.), tapuscrit (s.d.), copies d'articles (s.d.) — lecture modernisée du dossier entier}

\begin{document}

\section*{Résumé}

\begin{resume}
Ce dossier est un classeur de travail sur les courbes elliptiques --- les
courbes d'équation $y^2 = x^3 + ax + b$, qui sont aussi des groupes. Il
s'appelle « formulaire », et c'en est un : la plus grande partie des cent pages
n'est pas de Grothendieck. On y trouve, copiés ou dactylographiés, les tables
de John Tate (comment écrire l'équation, comment la changer de coordonnées,
comment reconnaître la « mauvaise réduction » d'une courbe modulo un nombre
premier), leur reprise par Pierre Deligne au-dessus d'une base quelconque,
deux articles d'Oliver Atkin sur des congruences entre coefficients de formes
modulaires découvertes à l'ordinateur, une liste des courbes à multiplication
complexe et une bibliographie. C'est la bibliothèque de quelqu'un qui se met
à un sujet.

Ce qui est de sa main tient en trois endroits, et c'est ce que cette lecture
restitue en entier.

D'abord, des questions au crayon dans les marges du texte de Deligne :
« pourquoi ? », « structure ? », et au bas de la dernière page cinq lignes qui
sont un programme --- l'invariant de Hasse, les formes modulaires sur un corps
quelconque, le relèvement canonique de Serre, la loi de groupe formel. C'est
une manière de lire : on lit un formulaire en lui posant des questions.

Ensuite, sept feuillets sur les \emph{périodes}. Une courbe elliptique
complexe est un tore $\mathbf{C}/\Lambda$, et sur un tore il y a deux façons de
compléter la différentielle $dz$ en une base des « formes à intégrer ». L'une
est algébrique : elle vient de l'équation, c'est $x\,dx/2y$. L'autre est
géométrique : c'est la conjuguée $d\bar z$. Une image : sur une feuille de
papier quadrillé, on peut choisir pour second axe celui que dessine la règle de
l'atelier, ou celui qui est perpendiculaire au premier ; ils diffèrent d'un
cisaillement, et ce cisaillement est un nombre. Les feuillets calculent ce
nombre $\alpha$, constatent qu'il ne dépend de rien d'arbitraire une fois
qu'on l'habille comme il faut, et demandent ce qu'il est. On sait aujourd'hui
le dire : c'est la série d'Eisenstein de poids deux \emph{corrigée} pour
devenir invariante, $G_2^*(\tau) = G_2(\tau) - \pi/\mathrm{Im}\,\tau$, qui n'est
plus holomorphe. Le calcul de la page est exact là où il s'arrête ; la
conclusion qu'il en tire ensuite sur la taille de $\alpha$ ne l'est pas, et
cette lecture dit pourquoi.

Enfin, cinq pages sur l'« espace » de toutes les courbes elliptiques à la fois
--- le \emph{site modulaire}, qu'on appelle aujourd'hui champ des modules. Deux
points y sont singuliers parce que leurs courbes ont plus de symétries que les
autres : $j=0$ et $j=12^3$. Au-dessus des nombres premiers $2$ et $3$, ces deux
points se confondent et les symétries deviennent $24$ ou $12$. Les pages
décrivent les deux lieux spéciaux, comptent les courbes qui y vivent sur
$\mathbf{Z}[\tfrac16]$ ($72$ et $32$), disent où ils cessent d'être lisses, puis
regardent ce qu'ils deviennent quand on marque les points d'ordre $3$ : quatre
branches qui se rejoignent toutes en un seul point au-dessus de $2$, et dont
l'argument montre qu'elles y ont la même tangente.

Les noms à chercher ensuite : algorithme de Tate, modèle de Néron, formes
modulaires entières de Deligne, décomposition de Hodge et quasi-périodes,
série d'Eisenstein non holomorphe, champ des modules $\mathcal{M}_{1,1}$,
invariant de Hasse, anneau de déformation universel, structures de niveau.

\keywords{Tate's algorithm, Kodaira symbols, minimal Weierstrass model, Néron model, integral modular forms, Atkin's congruences, p-adic Hecke operators, Hodge decomposition, de Rham cohomology of elliptic curves, quasi-periods, Legendre relation, non-holomorphic Eisenstein series, almost holomorphic modular forms, complex multiplication, class number one problem, moduli stack of elliptic curves, Hasse invariant, supersingular elliptic curves, universal deformation ring, level-3 structure, Picard group of the moduli stack}
\end{resume}

\pagerange{1}{100}
\section*{Le fil du dossier, et ses conventions}

Aucune page n'est numérotée par lui ; l'ordre est celui des feuillets, et il
n'est pas celui d'une argumentation. Le dossier se lit comme une pile.

\begin{enumerate}
\item \emph{Pages 2 à 13} --- copie de la lettre de J.~Tate à Cassels
(l'« algorithme de Tate »). Texte d'un autre, décrit.
\item \emph{Pages 15 à 33} --- le tapuscrit de P.~Deligne, « Courbes
elliptiques : formulaire, d'après J.~Tate », avec des marques au crayon dont
quelques-unes sont des questions rédigées. Décrit ; les questions sont lues.
\item \emph{Pages 35 à 39} --- cinq feuillets manuscrits de Tate, dans une
notation plus ancienne. Décrits.
\item \emph{Pages 40 à 63} --- deux tapuscrits d'A.~O.~L.~Atkin. Décrits.
\item \emph{Pages 65 à 76} --- \textbf{ses notes} : périodes, quasi-périodes,
et comparaison des deux décompositions de $H^1_{\mathrm{dR}}$. Sur le verso de
la page 74, un feuillet de l'avant-propos de SGA~4 corrigé de sa main (page
75).
\item \emph{Pages 79 et 80} --- liste dactylographiée des courbes à
multiplication complexe d'invariant rationnel, annotée au crayon.
\item \emph{Pages 82 à 93} --- bibliographie et « traduction libre d'un
manuscrit inédit de Tate ». Décrites.
\item \emph{Pages 95 à 99} --- \textbf{ses notes} : le site modulaire au
voisinage de $j=0$ et $j=12^3$, puis en niveau $3$.
\end{enumerate}

Les pages 1, 14, 34, 55, 64, 78, 81, 85, 94 sont blanches ; 67, 69, 71, 73, 77
et 100 sont des versos dactylographiés sans rapport (tirages Bourbaki, avant-propos
de SGA~4 barré, une page sur la spécialisation des cycles) ; 96 est une lettre
qui lui est adressée.

Les deux ensembles de notes ne se citent pas l'un l'autre, et ce n'est pas
cette lecture qui les reliera. Le premier est analytique (un seul tore
complexe, ses périodes) ; le second est arithmétique (toutes les courbes à la
fois, au-dessus de $\mathbf{Z}$). Leur point commun est le fibré $\omega$ des
différentielles invariantes, dont les puissances portent les formes modulaires :
$c_4$, $c_6$, $\Delta$ sont des sections de $\omega^{\otimes 4}$,
$\omega^{\otimes 6}$, $\omega^{\otimes 12}$ à la page 95, et l'invariant de la
page 66 est une section --- non holomorphe --- de $\omega^{\otimes 2}$.

\emph{Conventions, valables pour tout le dossier.} $E$ est une courbe
elliptique, $L = \underline{\omega}$ la droite de ses différentielles
invariantes (son $\underline{\omega}$), $\omega$ une base de $L$. Les poids
suivent la convention aujourd'hui usuelle (Katz) : une quantité de poids $k$
attachée au couple $(E,\omega)$ est multipliée par $\lambda^{-k}$ quand on
remplace $\omega$ par $\lambda\omega$, et devient ainsi une section
$f\,\omega^{\otimes k}$ indépendante du choix. Sur $\mathbf{C}$, $E =
\mathbf{C}/\Lambda$ avec $\Lambda = \mathbf{Z}\omega_1 + \mathbf{Z}\omega_2$,
$\tau = \omega_2/\omega_1$, $\mathrm{Im}\,\tau > 0$, et $\omega = dz$ ;
$V = H^1_{\mathrm{dR}}(E/\mathbf{C})$, extension de $L^{-1}$ par $L$. L'invariant
$j$ est partout l'invariant usuel $c_4^3/\Delta$, qui vaut $12^3$ pour la courbe
à quatre automorphismes --- non le $j = 1 - J$ du dossier 62, pages 18 à 53. Le
« site modulaire » $M$ est lu comme le champ des courbes de Weierstrass sans
fibre additive (le champ $\overline{\mathcal{M}}_{1,1}$ compactifié), puisque
la page 95 lui donne un diviseur à l'infini $M^{\infty} = \operatorname{div}
\Delta$.

Deux homonymies sont levées ici une fois pour toutes. Aux pages 72 à 76 il
réutilise les lettres $\alpha$, $\beta$ de la page 66 pour des objets voisins
mais distincts ; on les note ici $A$ et $B$. À la page 66 il note $\pi$ une
section canonique de $\underline{\omega}^{\otimes 2}$ ; on la note $\varpi$,
$\pi$ étant à la fois la différentielle invariante de Deligne et le nombre.

\pagerange{2}{63}
\section*{I. Les textes de référence (pages 2 à 63)}

\pagerange{2}{13}
\subsection*{La lettre de Tate (pages 2 à 13)}

Copie d'un manuscrit anglais de Tate --- table des matières et onze feuillets
--- qui est la lettre à Cassels publiée, légèrement retouchée, sous le titre
« Algorithm for determining the type of a singular fiber in an elliptic
pencil » (\emph{Modular Functions of One Variable IV}, LNM~476, 1975). On y
trouve l'équation de Weierstrass généralisée
\[
y^2 + a_1xy + a_3y = x^3 + a_2x^2 + a_4x + a_6,
\]
les quantités $b_2, b_4, b_6, b_8, c_4, c_6, \Delta$, $j = c_4^3/\Delta$ et les
relations $4b_8 = b_2b_6 - b_4^2$, $1728\,\Delta = c_4^3 - c_6^2$ ; les
changements de coordonnées $x = u^2x' + r$, $y = u^3y' + su^2x' + t$ et
l'invariance $u^{12}\Delta' = \Delta$ ; les équations minimales sur un anneau
de valuation discrète ; la filtration $E(K) \supset E_0(K) \supset E_1(K)
\supset \cdots$ et le groupe formel ; le facteur local $L_v(s)$ et la formule
$\int_{E(K)}|\omega| = (E(K):E_0(K))/L_v(1)$ ; le modèle de Néron et les dix
types de Kodaira ; enfin l'algorithme qui, en dix étapes, lit le type de la
fibre spéciale, le conducteur et l'indice $(E(K):E_0(K))$ sur les valuations des
coefficients.\footnote{Les marques en marge (« subprocedure », « sauf
erreur ») sont de la même main que le texte : rien sur ces feuillets n'est de
la sienne.}

\pagerange{15}{33}
\subsection*{Le formulaire de Deligne, et les questions dans ses marges (pages 15 à 33)}

Le tapuscrit, dix-neuf pages, est celui que Deligne a publié dans le même
volume sous le titre « Courbes elliptiques : formulaire d'après J.~Tate ». Il
reprend les formules de Tate au-dessus d'une base $S$ : la courbe $p\colon E
\to S$ de section unité $e$, le faisceau $\omega = e^*\Omega^1_{E/S}$, les
$p_*\mathcal{O}(ne)$ localement libres de rang $n$, l'équation obtenue par une
base adaptée de $p_*\mathcal{O}(3e)$ ; les formes réduites quand $2$ et $3$,
ou $2$ et $c_4$, sont inversibles, quand $2$ est nilpotent, en caractéristique
$2$ et $3$, avec les schémas de modules des courbes munies d'une
différentielle invariante, $\operatorname{Spec}\mathbf{Z}[\tfrac16][g_2,g_3]$
quand $2$ et $3$ sont inversibles (Proposition 2.5) et son analogue quand
seuls $2$ et $c_4$ le sont (3.5) ; puis les applications : la caractérisation des trois types de
fibres (Proposition 6.1), le schéma $\mathrm{Isom}_S(E,E')$ (6.3, 6.4), les
groupes d'automorphismes (6.9) ; enfin l'anneau des formes modulaires entières,
engendré par $c_4$, $c_6$, $\Delta$ sous la seule relation $c_4^3 - c_6^2 =
1728\,\Delta$ (7.1), et ses réductions $\mathbf{F}_2[a_1,\Delta]$ et
$\mathbf{F}_3[b_2,\Delta]$ (7.2).

Les marques au crayon sont pour la plupart des corrections d'épreuve et des
factorisations ($2^3 3^3$ sur $216$, $12^3$ sous $1728$). Quelques-unes sont
des phrases, d'une écriture rapide qui paraît être la sienne.\footnote{La
transcription laisse l'attribution ouverte, et nous de même : rien ne permet de
l'affirmer, et les factorisations, en particulier, pourraient être de
n'importe quel lecteur.} Lues dans l'ordre, elles disent ceci.

\emph{Le résumé des trois types} (page 20). Sur un corps, la fibre est
elliptique si $\Delta \neq 0$, multiplicative si $\Delta = 0$ et $c_4 \neq 0$,
additive si $\Delta = c_4 = 0$ ; et sous $\Delta = 0$, $c_4 = 0$ équivaut à
$c_6 = 0$. Cette dernière équivalence est vraie en toute
caractéristique,\footnote{En caractéristique $\neq 2,3$ elle découle de $c_4^3
= c_6^2$ ; en caractéristique $2$ et $3$, où la relation $1728\,\Delta = c_4^3 -
c_6^2$ est vide, elle découle des formules $c_4 = a_1^4$, $c_6 = a_1^6$
(caractéristique 2) et $c_4 = b_2^2$, $c_6 = -b_2^3$ (caractéristique 3) que le
tapuscrit donne en 7.2.} et le résumé est exact sur un corps quelconque, le
type étant géométrique.

\emph{Les valeurs spéciales de $j$} (page 26) : $j = 0 \Leftrightarrow c_4 = 0$
et $j = 12^3 \Leftrightarrow c_6 = 0$, là où $\Delta$ est inversible, puisque
$j = c_4^3/\Delta$ et $j - 12^3 = c_6^2/\Delta$. Au même endroit, « net » est
écrit au-dessus de « non ramifié » : c'est le mot des EGA pour la même notion.

\emph{Les automorphismes} (page 29). En face de l'énoncé selon lequel, $c_4$ et
$c_6$ étant inversibles, $E$ n'a que l'identité et la symétrie pour
automorphismes, la marge porte --- lecture incertaine --- « faux si $S$ non
connexe ». C'est juste si l'énoncé porte sur le groupe des automorphismes
au-dessus de $S$ : ce groupe est alors $(\mathbf{Z}/2)^{\pi_0(S)}$, et seul le
schéma en groupes $\underline{\mathrm{Aut}}_S(E)$ est $\mathbf{Z}/2$. En face du
groupe d'ordre $24$ de la caractéristique $2$, « structure ? » : ce groupe est
$\mathrm{SL}(2,\mathbf{F}_3)$, produit semi-direct du groupe des quaternions
$Q_8$ par $\mathbf{Z}/3$ ; celui d'ordre $12$ de la caractéristique $3$ est
$\mathbf{Z}/3 \rtimes \mathbf{Z}/4$, le facteur $\mathbf{Z}/4$ agissant par
inversion à travers son quotient d'ordre $2$.\footnote{Ces identifications sont
nôtres ; la page pose la question sans y répondre. Le groupe d'ordre $24$
reparaît, nommé $G'$, aux pages 97 à 99.}

\emph{Pourquoi ?} (page 31). En face de « il reviendrait au même de se limiter
aux courbes sans fibres additives ». La raison, que le tapuscrit ne donne pas
à cet endroit, est de codimension : dans l'espace $\operatorname{Spec}
\mathbf{Z}[a_1,\ldots,a_6]$ des équations, qui est normal, le lieu additif $c_4
= c_6 = 0$ est de codimension $2$, et une fonction définie hors de lui s'étend
(lemme de Hartogs algébrique).\footnote{Explication nôtre, proposée comme
réponse à la question de la marge.}

\emph{Le N.B. encadré} (page 31) : « sur le site modulaire, les diviseurs de
$c_4$, $c_6$, $\Delta$ sont à multiplicité $1$ ; celui de $\Delta$ est lisse sur
$\mathbf{Z}$, ceux de $c_4$, $c_6$ le sont sauf au-dessus des points $2$ et
$3$. » C'est le programme exact des pages 95 à 99, où il est démontré ; voir
la section~III. À la page 32, sous la Proposition 7.2 : le diviseur de $a_1$
est lisse sur $\mathbf{F}_2$, celui de $b_2$ sur $\mathbf{F}_3$.\footnote{La
page écrit $\mathbf{F}_2$ dans les deux lignes ; la seconde porte sur la
caractéristique $3$.}

\emph{Le programme} (page 33). Cinq lignes au bas de la dernière page :
\begin{itemize}
\item \emph{Interprétation de $a_1$, $b_2$ comme forme de Hasse-Witt ?} --- Oui :
la réduction modulo $2$ de $a_1$ et la réduction modulo $3$ de $b_2$ sont
l'invariant de Hasse, section de $\omega^{\otimes(p-1)}$ en caractéristique
$p$. En caractéristique $3$ c'est un calcul direct (le coefficient de $x^2$
dans le second membre de $y^2 = x^3 + \tfrac{b_2}{4}x^2 + \cdots$) ; en
caractéristique $2$, $a_1$ s'annule exactement sur le lieu supersingulier
$j=0$, à l'ordre $1$.
\item \emph{Formes modulaires sur un corps quelconque ? sur un anneau de
valuation discrète\ldots}, et \emph{Hasse-Witt sur un corps de caractéristique
$p>0$ quelconque ?} --- c'est la théorie des formes modulaires de Katz
(1973), où les formes modulaires sur un anneau $R$ quelconque sont définies
comme des règles $(E/R,\omega) \mapsto f(E,\omega) \in R$.
\item \emph{Relèvement canonique de Serre par formules universelles} --- le
relèvement canonique d'une courbe ordinaire (Serre--Tate).
\item \emph{Loi de groupe formel sur} (un symbole illisible)\emph{ : formules
explicites\ldots} --- la loi que la lettre de Tate écrit au §4.\footnote{Les noms donnés après chaque question sont les nôtres. Katz,
« $p$-adic properties of modular schemes and modular forms », LNM~350, est de
1973 ; Serre et Tate ont exposé le relèvement canonique à Woods Hole en 1964.
Rien sur la page ne dit qu'il connaissait l'un ou l'autre ; le symbole après
« sur » est illisible, et nous ne le restituons pas.}
\end{itemize}

\pagerange{35}{39}
\subsection*{Les feuillets de Tate en notation ancienne (pages 35 à 39)}

Cinq feuillets en anglais, de la main de Tate, où les $b_i$, $c_i$ s'appellent
encore $\beta_i$, $\gamma_i$ et le facteur d'échelle $\rho$ : l'équation
générale et ses transformations, puis les « formes canoniques spéciales » pour
$p \neq 2,3$, $p = 3$ et $p = 2$, avec les extensions qui rendent isomorphes
deux courbes de même $j$ et les groupes d'automorphismes, sous-groupes d'un
groupe d'ordre $24$. C'est le texte dont les pages 86 à 93 donnent une
traduction (section~II).

\pagerange{40}{63}
\subsection*{Les deux textes d'Atkin (pages 40 à 63)}

Deux tapuscrits d'A.~O.~L.~Atkin, en copie, avec les corrections d'un
correcteur d'épreuves et les instructions à l'imprimeur, qui font partie de
l'image photocopiée.\footnote{Aucune marque n'y est de sa main. La seule qui
porte un jugement mathématique, « multiplicat.\ semi group! » (page 58), est de
la main des instructions à l'imprimeur ; seul le « (1968) » à l'encre de la
page 40 a été ajouté sur la copie.} Le premier, « Congruences for modular
forms » (pages 40 à 54, daté 1968 à la main), raconte la recherche à
l'ordinateur Atlas de congruences pour les coefficients $c(n)$ de $j$ et pour
la fonction de partition $p(n)$ modulo les puissances d'un premier $q$ :
congruences « de Ramanujan » et congruences « multiplicatives », qui font
attendre une théorie de Hecke « $\mathfrak{p}$-adique ». Le second, « Congruence
Hecke operators » (pages 56 à 63), formule cette attente en conjectures pour le
groupe modulaire et pour $\Gamma_0(q)$, calquées sur les théorèmes de
Hecke--Petersson (son Théorème 1, base de formes propres) et sur leur
extension à $\Gamma_0(q)$ avec l'opérateur $U_q$ (Théorème 2), et espère, dans son dernier paragraphe,
une preuve par la géométrie algébrique.\footnote{C'est ce qu'ont fourni
ensuite les formes modulaires $p$-adiques de Serre et de Katz (1973), et l'on
comprend que ces deux textes aient été classés ici ; mais le dossier ne dit pas
ce qu'il en a tiré.}

\pagerange{65}{76}
\section*{II. Périodes et les deux décompositions de $H^1$ (pages 65 à 76)}

Sept feuillets de sa main, à l'encre, écrits au dos de feuillets
dactylographiés hors d'usage --- dont un projet barré de l'avant-propos de
SGA~4.\footnote{C'est un terme \emph{post quem} grossier pour ces notes (le
projet date au plus tôt de 1963--1964), non une date.}

\pagerange{65}{65}
\subsection*{Normalisation par une différentielle (page 65)}

Soit $E$ d'équation $y^2 = x^3 + \gamma_4 x + \gamma_6$ sur $\mathbf{C}$ et
$\omega = dx/2y$ sa différentielle invariante. Le couple $(E,\omega)$
détermine $x$ et $y$ ; si l'on remplace $\omega$ par $\lambda\omega$, on
remplace $x$, $y$, $\gamma_4$, $\gamma_6$ par $\lambda^{-2}x$, $\lambda^{-3}y$,
$\lambda^{-4}\gamma_4$, $\lambda^{-6}\gamma_6$. Les quantités
\[
x_2 = x\,\omega^{\otimes 2},\quad y_3 = y\,\omega^{\otimes 3},\quad
g_4 = \gamma_4\,\omega^{\otimes 4},\quad g_6 = \gamma_6\,\omega^{\otimes 6}
\]
ne dépendent donc pas de $\omega$, et vérifient $y_3^2 = x_2^3 + g_4x_2 +
g_6$ : c'est l'équation écrite en sections des puissances de
$L$.\footnote{La page écrit $x = \omega^{-2}x_2$, etc., en traitant $\omega$
comme un scalaire, et $dx = 2y\,\omega$ en le traitant comme une forme ; la
lecture par les poids est la nôtre, et c'est elle qui donne un sens aux deux
écritures à la fois. Plusieurs coefficients numériques sont biffés ou noircis
dans les deux cadres ; les formules se lisent comme données ici.} De $2y\,dy =
(3x^2+\gamma_4)\,dx$ on tire $dy = (3x^2+\gamma_4)\,\omega$. On pose
\[
\omega = \frac{dx}{2y}, \qquad \eta = x\,\omega = \frac{x\,dx}{2y},
\]
une forme de seconde espèce ($\eta$ a un pôle double sans résidu en
l'origine), de sorte que $(\omega,\eta)$ est une base de $V =
H^1_{\mathrm{dR}}(E/\mathbf{C})$.

Analytiquement, $E \simeq \mathbf{C}/\Lambda$ avec $\omega = dz$, $x =
\wp(z)$, $y = \tfrac12\wp'(z)$ :\footnote{La page écrit $y = -1/z^3 - \sum_{a} 1/(z-a)^3$ sans exclure $a=0$
de la somme, ce qui compterait deux fois le terme en $0$ ; le signe devant la
somme est surchargé et lu $-$, ce qui est le bon signe.}
\[
x = \frac{1}{z^2} + \sum_{a\in\Lambda\smallsetminus\{0\}}
\Bigl[\frac{1}{(z-a)^2} - \frac{1}{a^2}\Bigr], \qquad
y = -\sum_{a\in\Lambda}\frac{1}{(z-a)^3},
\]
et les périodes et quasi-périodes sont
\[
\omega_i = \int_{\gamma_i}\omega, \qquad
\eta_i = \int_{\gamma_i}\eta = \int_u^{u+\omega_i}\wp(z)\,dz
= -\bigl(\zeta(u+\omega_i) - \zeta(u)\bigr),
\]
indépendantes de $u$. Les $\eta_i$ de la page sont donc les opposées des
quasi-périodes de Weierstrass.\footnote{Avec cette convention, la relation de
Legendre s'écrit $\omega_1\eta_2 - \omega_2\eta_1 = 2\pi i$ pour
$\mathrm{Im}(\omega_2/\omega_1) > 0$. La page ne la cite pas ; elle sert
plus bas. Les seconds membres de « $\gamma_4 = $ », « $\gamma_6 = $ » sont
restés en blanc.}

\pagerange{66}{66}
\subsection*{Le passage du scindage algébrique au scindage de Hodge (page 66)}

$V$ contient la droite $L = \mathbf{C}\omega$ (la filtration de Hodge) et
possède deux supplémentaires naturels de $L$ : la droite $\mathbf{C}\eta$, qui
vient de l'équation, et la droite $\mathbf{C}\bar\omega = H^{0,1}$, qui vient
de la conjugaison complexe. On écrit donc
\[
\eta + \alpha\,\omega = \beta\,\bar\omega .
\]
En intégrant sur $\gamma_1$ et $\gamma_2$ et en éliminant $\beta$ :
\[
\alpha = -\frac{\bar\omega_1\eta_2 - \bar\omega_2\eta_1}
{\bar\omega_1\omega_2 - \bar\omega_2\omega_1},
\qquad
\bar\omega_1\omega_2 - \bar\omega_2\omega_1 = 2i\,|\omega_1|^2\,\mathrm{Im}\,\tau
\neq 0 .
\]
Remplacer $\omega$ par $\lambda\omega$ remplace $\eta$ par $\lambda^{-1}\eta$ et
$\alpha$ par $\lambda^{-2}\alpha$ : la section
\[
\varpi = \alpha\,\omega^{\otimes 2} \in \Gamma(\underline{\omega}^{\otimes 2})
= \Gamma\,\mathrm{Hom}(\underline{\omega}^{-1},\underline{\omega})
\]
est canonique.\footnote{La page la note $\pi$.} C'est la différence des deux
supplémentaires : les supplémentaires de $L$ dans $V$ forment un espace affine
sous $\mathrm{Hom}(V/L, L) = \mathrm{Hom}(L^{-1},L) = L^{\otimes 2}$, et
$\varpi$ est le vecteur qui va du premier au second.\footnote{L'identification
$V/L \simeq L^{-1}$ est celle de la dualité : $\eta$ y est la base duale de
$\omega$ pour l'accouplement $\tfrac{1}{2\pi i}\int_E$, par la relation de
Legendre. Cette lecture de $\varpi$ comme différence de scindages est la
nôtre ; la page la dit à la page 74 (l'élément $\gamma$), dans le langage des
pages 72 à 74.}

Deux questions suivent : \emph{a) Quand $\alpha = 0$, i.e.\ $\bar\omega$
proportionnel à $\eta$ ? b) Quand $\alpha$ est-il algébrique sur
$\mathbf{Q}$ ?}

\emph{Ce qu'est $\alpha$.} La relation de Legendre donne la forme close
qu'aucune des pages n'atteint. Si $\eta'_i = \zeta(z+\omega_i) - \zeta(z)$,
on a $\eta'_1 = \omega_1^{-1}G_2(\tau)$ avec $G_2(\tau) = \sum_n
\sum_{m}{}' (m+n\tau)^{-2}$, et $\eta'_1\omega_2 - \eta'_2\omega_1 = 2\pi i$ ;
en reportant dans la formule de $\alpha$,
\[
\alpha = \frac{1}{\omega_1^2}\Bigl(G_2(\tau) - \frac{\pi}{\mathrm{Im}\,\tau}\Bigr)
= \frac{G_2^*(\tau)}{\omega_1^2},
\qquad
\beta = -\frac{\pi}{|\omega_1|^2\,\mathrm{Im}\,\tau}.
\]
La section canonique $\varpi$ est donc la série d'Eisenstein non holomorphe de
poids $2$, $G_2^* = \tfrac{\pi^2}{3}E_2^*$ : la seule correction de $G_2$ qui
soit modulaire, au prix de l'holomorphie.\footnote{Calcul nôtre ; la page
s'en tient aux séries des pages 68 et 70. La correction $-\pi/\mathrm{Im}\,\tau$
remonte à Hecke (1927). L'interprétation de $E_2^*$ par le scindage de Hodge,
et de $E_2$ par le scindage de la racine unité dans le cas $p$-adique, est
celle de Katz (1973, 1976) ; ces pages en ont la version complexe, sans le
nom. Nous ne disons pas qu'il l'ait reçue de Katz, ni l'inverse.} Les deux
questions reçoivent alors une réponse partielle. a) $\alpha$ s'annule aux
points $\tau = i$ et $\tau = e^{2\pi i/3}$ et à leurs transformés par
$\mathrm{SL}_2(\mathbf{Z})$, car $G_2^*$ est modulaire de poids $2$ et ces
points sont fixes par des éléments d'ordre $4$ et $6$ ;\footnote{Par exemple
$G_2^*(-1/\tau) = \tau^2G_2^*(\tau)$ donne en $\tau = i$ : $G_2^*(i) =
-G_2^*(i)$. Nous ne disons pas s'il y a d'autres zéros.} en ces points
$\bar\omega$ est bien proportionnel à $\eta$. b) Si $E$ est définie sur
$\overline{\mathbf{Q}}$ et a multiplication complexe, et si $\omega$ est
algébrique, $\alpha$ est algébrique : les endomorphismes décomposent $V$ en
deux droites propres définies sur $\overline{\mathbf{Q}}$, dont l'une est $L$
et l'autre est nécessairement $H^{0,1}$.\footnote{La réciproque n'est pas
abordée ici.}

\pagerange{68}{70}
\subsection*{Les quasi-périodes en séries (pages 68 et 70)}

La page 68 calcule $\eta_1$ en primitivant terme à terme, $\int x\,dz = -1/z
+ \sum'\bigl[-1/(z-a) - z/a^2\bigr]$, entre $u$ et $u+\omega_1$. Le terme
général devient $\omega_1\bigl[\tfrac{1}{(u-a)(u+\omega_1-a)} -
\tfrac{1}{a^2}\bigr]$, qui est $O(|a|^{-3})$ : la série converge absolument,
et l'on peut faire $u \to 0$ après avoir isolé les termes $a = 0$ et $a =
\omega_1$, qui donnent ensemble $-3/\omega_1$. On obtient
\[
\eta_1 = -\frac{3}{\omega_1} + \omega_1^2
\sum_{a\neq 0,\ \omega_1}\frac{1}{a^2(a-\omega_1)}, \qquad
\eta_2 = -\frac{3}{\omega_2} + \omega_2^2
\sum_{a\neq 0,\ \omega_2}\frac{1}{a^2(a-\omega_2)},
\]
avec $a = m\omega_1 + n\omega_2$.\footnote{Pour $\eta_2$, l'indice de la
somme se lit « $a \neq \omega_1$ », recopié de la ligne précédente ; c'est
$a \neq \omega_2$ qu'il faut. Les deux formules sont exactes : nous les avons
vérifiées.} Ces séries, contrairement à celle qui définit $G_2$, convergent
absolument ; c'est l'avantage du chemin pris.

La page 70 forme $\bar\omega_1\eta_2 - \bar\omega_2\eta_1$ sur l'ensemble
d'indices commun $(m,n) \neq (0,0), (1,0), (0,1)$, en mettant à part les deux
termes que les deux sommes ne partagent pas, puis passe en $\tau$. La partie
constante est $-3(\bar\omega_1/\omega_2 - \bar\omega_2/\omega_1) =
-3\,(|\omega_1|^2 - |\omega_2|^2)/\omega_1\omega_2$, et les termes isolés
valent $(\bar\omega_1\omega_2^4 + \bar\omega_2\omega_1^4)/\omega_1^2\omega_2^2
(\omega_1-\omega_2)$. La page conduit le calcul jusqu'à une expression de
$\alpha$ sous la forme $-\bigl(2i\,\omega_1^2\,\mathrm{Im}\,\tau\bigr)^{-1}$
fois une fonction de $\tau$ seule --- ce qui exhibe le poids $2$ ---, sans
sommer la série.\footnote{À partir de la troisième ligne, le facteur $3$ de
$-3(1-|\tau|^2)/\tau$ disparaît ; la forme close de la page 66 ne dépend pas
de ce détail, et nous ne reprenons pas la série terme à terme, dont la page
note elle-même plusieurs états surchargés (signes, exposants).}

\pagerange{72}{74}
\subsection*{L'algèbre linéaire des deux scindages (pages 72 à 74)}

Les pages 72 et 74 reprennent la situation sans coordonnées. $V$ est un espace
de dimension $2$ muni d'une forme alternée, extension $0 \to L \to V \to
L^{-1} \to 0$, et la conjugaison complexe est une involution antilinéaire,
c'est-à-dire un isomorphisme $\varphi\colon \bar V \to V$ avec $\varphi\bar\varphi
= 1$. Une fois choisi un scindage algébrique $V = L \oplus L^{-1}$, la
restriction de $\varphi$ à $\bar L$ a deux composantes,
\[
A\colon \bar L \to L, \quad A \in \bar L^{-1}\otimes L ; \qquad
B\colon \bar L \to L^{-1}, \quad B \in (\bar L L)^{-1}.
\]
\footnote{La page les nomme $\alpha$ et $\beta$, comme à la page 66 ; elle
indique sous $\bar L^{-1}\otimes L$ « défini par $T^{(2)}$ sur $\mathcal{U}$ »
et sous $\bar L L$ « défini par $P^{(2)}$ sur $\mathbf{R}^{*+}$ », $L$ étant
« défini par $T$ sur $\mathcal{U}$, $P$ sur $\mathbf{R}^{*+}$ ». Nous lisons
la décomposition polaire $\mathbf{C}^* = \mathcal{U} \times \mathbf{R}^{*+}$ :
une droite complexe est donnée par un torseur $T$ sous le cercle
$\mathcal{U}$ et un torseur $P$ sous $\mathbf{R}^{*+}$ ; $\bar L^{-1} \otimes
L$ ne dépend que de $T^{\otimes 2}$ (son module est canoniquement trivial) et
$\bar L L$ que de $P^{\otimes 2}$ (son argument l'est). Cette lecture est la
nôtre ; plusieurs mots de la ligne sont illisibles.}
L'hypothèse $B \neq 0$ (le scindage algébrique n'est pas celui de Hodge, ce
qui est toujours le cas puisque $\bar L \cap L = 0$) permet d'exprimer
$\varphi$ sur $L^{-1}$ : si $y \in L^{-1}$, $y = -AB^{-1}(y) +
\varphi\bigl(\overline{B^{-1}(y)}\bigr)$. La page demande alors quand la
conjugaison respecte la forme alternée, $\langle \bar x, \bar y\rangle =
\overline{\langle x, y\rangle}$, et trouve, en l'écrivant pour $x \in L$, $y
\in L^{-1}$, la condition $B + \bar B = 0$ : \emph{$B$ est purement
imaginaire}.\footnote{Avec la normalisation $\langle\omega,\eta\rangle = 1$ ;
dans la page, les deux lignes qui précèdent cette conclusion sont très
surchargées et nous ne reprenons que leur état final, qui suffit.}

Dans le cas d'une courbe elliptique (page 74), $B$ est la forme
$\omega \mapsto \int_E \omega\wedge\bar\omega$, et pour $\omega = dz =
ds + i\,dt$ ($s$, $t$ réels),
\[
\int_E \omega\wedge\bar\omega = -2i\int_E ds\wedge dt
= -2i\,|\omega_1|^2\,\mathrm{Im}\,\tau,
\]
donc
\[
B = -2i\rho,\qquad \rho = \int_E ds\wedge dt = \mathrm{Vol}(E) > 0 .
\]
\footnote{La page écrit $dx + i\,dy$, avec $x = \mathrm{Re}\,z$, $y =
\mathrm{Im}\,z$ ; nous changeons de lettres pour ne pas confondre avec les
coordonnées de l'équation.}
Quant à $A$, c'est \emph{a priori} un élément quelconque de $\bar L^{-1}\otimes
L$, défini indépendamment du choix de $\omega$ ; la page demande ce qu'il est.
Elle note que $\gamma = -AB^{-1} \in L^{\otimes 2}$ est l'objet qui fait passer
du scindage algébrique au scindage transcendant, et que la donnée de $(A,B)$
équivaut à celle de $B$ et $\gamma$. C'est la section $\varpi$ de la page 66,
au signe près.\footnote{Le supplémentaire de Hodge $\bar L$ est le graphe de
$AB^{-1}\colon L^{-1} \to L$ ; le signe de $\gamma$ dépend du sens dans
lequel on écrit le passage. En coordonnées, $A(\bar\omega) =
(\alpha/\beta)\,\omega$ avec les $\alpha$, $\beta$ de la page 66 ; d'où
$|A| = \mathrm{Im}\,\tau\,|G_2^*(\tau)|/\pi$, qui ne dépend que de $\tau$.}

\pagerange{74}{76}
\subsection*{Module et argument de $A$ (pages 74 et 76)}

La page pose, pour $\omega$ et $\eta$ choisis,
\[
\rho = -\frac{1}{2i}\int_E \omega\wedge\bar\omega = \mathrm{Vol}(E), \qquad
\sigma = -\frac{1}{2i}\int_E \eta_0\wedge\bar\eta_0,
\]
où $\eta_0$ est une forme fermée de la classe de $\eta$ ($\sigma$ ne dépend pas
de ce choix). Lorsque $\eta$ est normalisée par $\int_E \omega\wedge\eta = 1$,
les deux conditions $\varphi\bar\varphi = 1$ et $B + \bar B = 0$ donnent
\[
4\rho\sigma = A\bar A - 1 .
\]
\footnote{La page écrit $4\rho\sigma = 1 - \alpha\bar\alpha$, et en conclut
(page 76) $0 \leq \alpha\bar\alpha \leq 1$. Le signe est faux, et la
conclusion aussi : $\sigma$ n'a pas de signe. En $\tau = i$, $A = 0$ et $\sigma
= -1/4\rho < 0$ ; quand $\mathrm{Im}\,\tau \to \infty$, $G_2^* \to \pi^2/3$ et
$|A| \sim \pi\,\mathrm{Im}\,\tau/3 \to \infty$. Avec $\eta = x\,dx/2y$
elle-même, pour laquelle $\int_E\omega\wedge\eta = 2\pi i$, la même relation
devient $4\rho\sigma = 4\pi^2(A\bar A - 1)$ : le signe ne dépend pas de la
normalisation.} Ainsi $B$ et $A\bar A$ sont donnés par des volumes ; il reste
l'argument de $A$ (pour $A \neq 0$), qui dépend du choix de $\omega$.

La page normalise alors $\omega$ par $\omega^{\otimes 12} = \Delta$, ce qui le
détermine à une racine douzième de l'unité $\zeta$ près. Comme $A \in \bar
L^{-1}\otimes L$, son écriture numérique dans la base $\omega\otimes\bar
\omega^{-1}$ est multipliée par $\zeta/\bar\zeta = \zeta^2$ : elle est
déterminée à une racine sixième de l'unité près, et
\[
\Bigl(\frac{A}{|A|}\Bigr)^{6} \in \mathcal{U} = \{z \in \mathbf{C},\ |z| = 1\}
\]
est un invariant canonique de la courbe. La normalisation fixe aussi $\rho$ et
$\sigma$. La page demande comment l'exprimer en fonction de $j$, sans
répondre. Avec la forme close de la page 66, la réponse est
\[
\Bigl(\frac{A}{|A|}\Bigr)^{6} = \frac{G_2^*(\tau)^6/\Delta(\tau)}
{\bigl|G_2^*(\tau)^6/\Delta(\tau)\bigr|},
\]
fonction de $\tau$ invariante par $\mathrm{SL}_2(\mathbf{Z})$, donc fonction sur
la droite des $j$ --- mais analytique réelle, non holomorphe : c'est une
fonction de $j$ et de $\bar j$, non de $j$ seul.\footnote{Calcul nôtre :
$A = -(\bar\omega_1/\omega_1)\,\mathrm{Im}\,\tau\;G_2^*(\tau)/\pi$, et
$\omega^{\otimes 12} = \Delta$ revient à $\omega_1^{12} = \Delta(\tau)$. La
fin de la phrase de la page, « (\ldots{} comment \ldots{} en fonction de $j$ !) »,
est en partie illisible.}

\pagerange{75}{75}
\subsection*{Un feuillet de l'avant-propos de SGA 4 (page 75)}

Au dos de la page 74, un feuillet dactylographié barré du projet
d'avant-propos de SGA~4 (travaux d'Igusa, d'Ogg et Chafarevitch ; plan des
exposés I à XIX), qui porte ses corrections à l'encre : ponctuation, numéros
d'exposés repassés, « théorème des dualités globales » ramené à « théorème de
dualité globale », et, après « Théorèmes de changement de base », l'ajout
« qui techniquement constituent le résultat central de ce séminaire ». Rien de
mathématique ne s'y rapporte aux courbes elliptiques ; la transcription donne
le détail.

\pagerange{79}{80}
\section*{III. Les courbes à multiplication complexe d'invariant rationnel (pages 79 et 80)}

Un tapuscrit anonyme de deux feuillets, paginés « X-1 » et « X-2 », dresse la
liste des courbes à multiplication complexe dont $j$ est rationnel : elles
correspondent aux ordres de nombre de classes $1$ des corps quadratiques
imaginaires, $\mathbf{Q}(\sqrt{-d})$ pour $d = 1, 2, 3, 7, 11, 19, 43, 67,
163$ et les conducteurs $f = 2$ ($d = 1, 3, 7$) et $f = 3$ ($d = 3$) --- treize
ordres, et treize valeurs de $j$, de $j(i) = 12^3$ à $j\bigl(\tfrac{-1+\sqrt{-163}}{2}\bigr) =
-(2^6\cdot 3\cdot 5\cdot 23\cdot 29)^3$, tirées du traité de Weber. Le texte
réserve « peut-être une autre valeur $d_?$ ».\footnote{Heegner (1952), Baker
(1966) et Stark (1967) ont montré qu'il n'y en a pas. La réserve date le texte
d'avant ces résultats, ou d'un auteur qui ne les tenait pas pour acquis.}

Au crayon, deux remarques. Sur $d = 1, 2, 3, 7, 11$, « euclidiens » : ce sont
en effet exactement les corps quadratiques imaginaires euclidiens pour la
norme. En marge : « NB les $d$ sont premiers ! » --- vrai pour $d > 1$ ; la
raison est la théorie des genres, selon laquelle le $2$-rang du groupe des
classes est le nombre de premiers divisant le discriminant, moins un, de sorte
qu'un nombre de classes $1$ force un discriminant divisible par un seul
premier.\footnote{Le cas $d = 1$ (discriminant $-4$) échappe à la lettre de la
remarque, non à son sens. L'explication par les genres est nôtre.}

\pagerange{82}{93}
\section*{IV. Bibliographie et formulaire traduit (pages 82 à 93)}

Un tapuscrit anonyme en deux suites. La première (pages 82 à 84) est une
bibliographie en dix rubriques, des fonctions elliptiques classiques (Jordan,
Hurwitz--Courant, Fricke) à la multiplication complexe « formelle » de Lubin
et Tate (1965), en passant par les schémas modulaires d'Igusa et, parmi les
variétés abéliennes, deux exposés de ses \emph{Fondements de la géométrie
algébrique} (232 et 236) ; quelques traits au crayon en marge, sans mots. La
seconde (pages 86 à 93) se donne pour « Courbes elliptiques -- Formulaire
(traduction libre d'un manuscrit inédit de Tate) », et suit en effet, en
français, les feuillets de Tate des pages 35 à 39 : formules générales,
changements de coordonnées, formes canoniques pour $p \neq 2,3$, $p=3$, $p=2$,
« courbes types », groupes d'automorphismes d'ordre $12$ et $24$, anneaux
d'endomorphismes (un ordre maximal du corps de quaternions ramifié en $3$ et
$\infty$ ; les quaternions entiers de Hurwitz en caractéristique $2$).\footnote{Le
traducteur n'est pas nommé. La lettre de Tate renvoie (page 4) au « Courbes
Elliptiques -- Formulaire » de Serre ; il est tentant d'y reconnaître ce
tapuscrit, mais rien sur les feuillets ne l'établit. Les seules marques
manuscrites sont des factorisations ($2^2.3$ sous $12$, $2^6 3^3$ sous $1728$)
et, page 91, deux fins d'équation gribouillées qui font lire $r^3 + a_4r = 0$
et $r^3 + a_4r + 2a_6 = 0$.}

\pagerange{95}{99}
\section*{V. Le site modulaire près de $j = 0$ et $j = 12^3$ (pages 95 à 99)}

Cinq pages de sa main, de la même encre. Elles réalisent le N.B. de la page 31
(section~I). La page 96 est une lettre datée du 9 mars 1968 ; si les pages
97 et suivantes sont écrites au dos de cette feuille, elles lui sont
postérieures.\footnote{Conditionnel de la transcription, que nous gardons.}

$M$ désigne ici le champ des courbes de Weierstrass sans fibre additive au-dessus
de $\mathbf{Z}$ --- le champ des modules compactifié
$\overline{\mathcal{M}}_{1,1}$ ---, qu'il appelle « site modulaire ».\footnote{La
notion de champ algébrique est celle de Deligne et Mumford (1969) ; « site
modulaire » est son nom d'alors pour l'objet, vu par sa topologie. Le
compactifié est imposé par la page elle-même, qui parle du diviseur
$M^{\infty} = \operatorname{div}\Delta$ : sur le champ des seules courbes
elliptiques, $\Delta$ est inversible.}

\pagerange{95}{95}
\subsection*{Les trois diviseurs et leurs classes (page 95)}

Un dessin montre, au-dessus de $\operatorname{Spec}\mathbf{Z}$, le diviseur
$\Delta = 0$ (la pointe), le diviseur $j = 0$, c'est-à-dire $c_4 = 0$, et le
diviseur $j = 12^3$, c'est-à-dire $c_6 = 0$, qui se coupent au-dessus de $p =
2$ et de $p = 3$ --- puisque $12^3 = 2^6 3^3$ ---, en des points où le groupe
d'automorphismes est d'ordre $24$ et $12$. En caractéristique $\neq 2,3$, les
points de $j = 0$ ont $6$ automorphismes et ceux de $j = 12^3$ en ont $4$. Les
diviseurs se lisent :
\[
\operatorname{div}(j)^{+} = 3\operatorname{div}(c_4), \quad
\operatorname{div}(j-12^3)^{+} = 2\operatorname{div}(c_6), \quad
\operatorname{div}(j)^{-} = \operatorname{div}(j-12^3)^{-} = M^{\infty} =
\operatorname{div}\Delta,
\]
puisque $j = c_4^3/\Delta$ et $j - 12^3 = c_6^2/\Delta$ ; et leurs classes dans
le groupe de Picard du champ sont
\[
\mathrm{cl}\bigl(\operatorname{div} c_4\bigr) = \omega^{4}, \qquad
\mathrm{cl}\bigl(\operatorname{div} c_6\bigr) = \omega^{6}, \qquad
\mathrm{cl}\bigl(M^{\infty}\bigr) = \omega^{12},
\]
$c_4$, $c_6$, $\Delta$ étant des sections de $\omega^{\otimes 4}$,
$\omega^{\otimes 6}$, $\omega^{\otimes 12}$.\footnote{Le groupe de Picard du
champ non compactifié est $\mathbf{Z}/12$, engendré par $\omega$ (Mumford,
1965, sur un corps ; Fulton et Olsson, 2010, sur $\mathbf{Z}$). Rien sur la
page ne dit qu'il connaissait le résultat de Mumford.}

\emph{Le lieu $M' = V(c_4)$}, sous-champ fermé défini par $c_4 = 0$.

a) Au-dessus de $\mathbf{Z}[\tfrac16]$, $M'$ est isomorphe au champ classifiant
$B\mu_6$ : une courbe elliptique $E$ sur une base $S$ où $6$ est inversible,
avec $c_4(E) = 0$, est une forme tordue de $y^2 = x^3 + 1$, et lui correspond
le $\mu_6$-torseur de ses isomorphismes avec cette courbe, dont le groupe
d'automorphismes est $\mu_6$ agissant par $(x,y) \mapsto (\zeta^2x,\zeta^3y)$.
Il y a donc, sur $\mathbf{Z}[\tfrac16]$, autant de telles courbes que
d'éléments de
\[
H^1\bigl(\operatorname{Spec}\mathbf{Z}[\tfrac16], \mu_6\bigr)
\simeq \mathbf{Z}[\tfrac16]^{*}/\mathbf{Z}[\tfrac16]^{*6}
\simeq \mathbf{Z}/2 \times \mathbf{Z}/6 \times \mathbf{Z}/6,
\]
soit $72$ : les courbes $y^2 = x^3 + u$, $u$ parcourant les unités $\pm
2^a3^b$ modulo les puissances sixièmes.\footnote{Le groupe de Picard de
$\mathbf{Z}[\tfrac16]$ étant nul, la suite de Kummer donne le premier
isomorphisme ; les unités sont $\{\pm 1\}\times 2^{\mathbf{Z}}\times
3^{\mathbf{Z}}$. La page écrit $(\mathbf{Z}[\tfrac16])^*_6$ ; le crochet
qu'elle ouvre n'est pas refermé.}

b) La fibre de $M'$ en $p = 2$ est de multiplicité $4$ en un point unique,
celle en $p = 3$ de multiplicité $2$ : $M'$ est ramifié sur $\mathbf{Z}$
exactement en ces deux points. Il n'est pas régulier en $p = 2$, il l'est en $p
= 3$.\footnote{Vérification : en caractéristique $2$, $c_4 = a_1^4 + 8a_1^2a_2
+ 16a_2^2 - 24a_1a_3 - 48a_4$ appartient au carré de l'idéal $(2, a_1)$,
d'où la multiplicité $4$ et la non-régularité ; en caractéristique $3$, $c_4 =
b_2^2 - 24b_4$ où $b_4 = 2a_4$ est une unité au point supersingulier
($\Delta = -a_4^3 \neq 0$), de sorte que $c_4$ est un paramètre régulier.}

\emph{Le lieu $M'' = V(c_6)$.} a) Au-dessus de $\mathbf{Z}[\tfrac16]$, $M''
\simeq B\mu_4$, par la courbe $y^2 = x^3 - x$ ; il y a $32$ courbes elliptiques
sur $\mathbf{Z}[\tfrac16]$ avec $c_6 = 0$, les $y^2 = x^3 + ux$, $u$ modulo
les puissances quatrièmes ($|\mathbf{Z}/2\times\mathbf{Z}/4\times\mathbf{Z}/4|
= 32$).\footnote{La fin de la phrase de la page, « \ldots{} choisir comme
origine\ldots », est en partie illisible.} b) $M''$ est de multiplicité $6$ en
$p = 2$ et $3$ en $p = 3$, ramifié exactement en ces deux points, et régulier
en aucun des deux --- ce que donnent $c_6 \equiv a_1^6 \pmod 2$ et $c_6
\equiv -b_2^3 \pmod 3$, tous les autres termes de $c_6 = -b_2^3 + 36b_2b_4 -
216b_6$ étant dans le carré de l'idéal maximal.

\pagerange{98}{98}
\subsection*{Paramètres locaux (page 98)}

Au-dessus de $\mathbf{Z}$, $M$ est de dimension relative $1$, et la page donne
un paramètre local relatif en chaque région :
\begin{itemize}
\item hors de $M' \cup M'' \cup M^{\infty}$ : $j$ ;
\item hors de $M' \cup M''$, au voisinage de la pointe : $1/j$ ou $1/(j - 12^3)$ ;
\item au voisinage de $M'$ en caractéristique $\neq 2,3$ : $c_4$ ; de $M''$ : $c_6$ ;
\item au point $p = 2$, $c_4 = c_6 = 0$ : $a_1$ ; au point $p = 3$ : $b_2$.
\end{itemize}
Il demande ensuite la structure de l'anneau local complété de $M$ en un point
géométrique où $c_4 = c_6 = 0$, et propose $W[[t]]$, $W$ l'anneau des
vecteurs de Witt du corps résiduel ($\overline{\mathbf{F}}_2$ ou
$\overline{\mathbf{F}}_3$), avec $t = a_1$ ou $t = b_2$, en laissant à écrire
les équations de $c_4 = 0$ et $c_6 = 0$. C'est l'anneau de déformation
universel de la courbe supersingulière, et la réponse est exacte ; les
équations sont, à des unités et à des termes de $(p,t)^2$ près, $c_4 \equiv t^4$,
$c_6 \equiv t^6$ en caractéristique $2$ et $c_4 \equiv t^2$, $c_6 \equiv -t^3$
en caractéristique $3$, ce qui redonne les multiplicités de la page 95.\footnote{Le
nom et la référence sont nôtres (Lubin--Tate, 1966 ; Serre--Tate). Le
paramètre $t$ relève l'invariant de Hasse, ce qui rejoint la première question
de la page 33.}

\pagerange{98}{99}
\subsection*{Le niveau 3 au-dessus de $j = 0$ (pages 98, 97 et 99)}

Soit $M_3$ le champ --- ici un schéma --- des courbes elliptiques munies d'une
base $(e_1,e_2)$ de leurs points d'ordre $3$, au-dessus de $\mathbf{Z}[\tfrac13]$ ;
l'accouplement de Weil $\langle e_1, e_2\rangle$ le fait vivre au-dessus du
schéma $\mu_3^{*}$ des racines primitives cubiques de l'unité, et $G =
\mathrm{GL}(2,\mathbf{F}_3)$ y opère.\footnote{La page ne définit pas $M_3$ ;
c'est la lecture de la transcription, que tout le contexte confirme. Les
ordres sont ceux de la colonne en marge de la page 98 : $|G| = 48$, le Borel
$B$ d'ordre $12$, $G/B$ d'ordre $4$, $G' = \mathrm{SL}(2,\mathbf{F}_3)$ d'ordre
$24$, $B' = G'\cap B$ d'ordre $6$.}

\emph{La structure de $M_3 \cap V(c_4)$ hors de $6$} (page 98). Si $E$ est une
courbe elliptique sur $S$, $6$ inversible, avec $c_4(E) = 0$, son groupe
d'automorphismes $\mu_6 = \mu_2\times\mu_3$ agit sur ${}_3E$ par $\pm 1$ (le
facteur $\mu_2$, central) et par un élément unipotent (le facteur $\mu_3$, qui
est $B_u$) ; il y a donc dans ${}_3E$ un unique drapeau invariant, ce qui réduit
le groupe structural de $G$ au Borel $B$. Sur $\mathbb{S} =
\mu_3^{*}[\tfrac16]$ il existe une courbe $\mathbb{E}$ avec $c_4 = 0$ et une base
adaptée au drapeau ; toute courbe avec $c_4 = 0$ et structure de niveau $3$
s'en déduit localement, et
\[
M_3[\tfrac16] \cap V(c_4) \;\simeq\; \mu_3^{*}[\tfrac16] \overset{B}{\times} G,
\]
réunion disjointe de $|G/B| = 4$ copies de $\mu_3^{*}[\tfrac16]$.\footnote{Une
bonne partie des phrases qui relient ces énoncés sont illisibles ; nous
restituons l'enchaînement que les formules lisibles imposent. Le décompte se
vérifie directement : à accouplement fixé, les structures de niveau forment
un torseur sous $G'$, sur lequel $\mathrm{Aut}(E) = \mu_6$ agit librement
par $B'$ ; il y a $24/6 = 4$ orbites.}

\emph{En $p = 2$} (pages 97 et 99). La courbe $\mathbb{E}$ a bonne réduction en
$2$ sur l'anneau des entiers de $\mathbf{Q}(\sqrt[3]{1})$ (par exemple $y^2 + y
= x^3$, de discriminant $-27$).\footnote{La page ajoute « (ni 3 ?) » ; la
question est hors de propos pour $M_3$, qui ne vit qu'au-dessus de
$\mathbf{Z}[\tfrac13]$.} Les quatre composantes irréductibles de $M'_3 = M_3
\cap V(c_4)$ ne se rencontrent qu'au-dessus de $2$, et s'y rencontrent toutes
en un seul point $a$, au-dessus du point $b = \operatorname{Spec}\mathbf{F}_4$
de $\mu_3^{*}$ : la courbe supersingulière de caractéristique $2$ a pour
groupe d'automorphismes $G' = \mathrm{SL}(2,\mathbf{F}_3)$ tout entier, qui
identifie entre elles toutes ses structures de niveau d'accouplement donné.
$M'_3$ est connexe.

Reste à savoir si les quatre branches ont en $a$ des tangentes distinctes.
L'espace tangent de Zariski de $M_3$ en $a$ est de dimension $2$ sur
$\mathbf{F}_4$ et se scinde en $V_1 \times V_2$, $V_1$ vertical (tangent à la
fibre $(M_3)_b$) et $V_2$ horizontal. Une branche, qui est une section de
$M_3 \to \mu_3^*$, a une tangente non verticale, c'est-à-dire le graphe d'une
application linéaire $V_2 \to V_1$ : il y a exactement $4$ telles droites,
$\operatorname{card}\mathbf{P}^1(\mathbf{F}_4) - 1 = 5 - 1$. Or $G'$ opère
trivialement sur $V_2$ (il respecte l'accouplement, donc la base
$\mu_3^{*}$) et sur $V_1$ par un caractère $G' \to \mathbf{F}_4^{*}$ ; le noyau
en est le $2$-sous-groupe de Sylow $P' \simeq Q_8$, d'indice $3$, de sorte que
$G' \to \mathbf{F}_4^{*} \simeq \mathbf{Z}/3$ est surjectif.\footnote{La page
écrit « le $2$-Sylow $P'$ de $G'$, qui est d'indice 8 » puis
« $\operatorname{card}[G':P'] = 3$ » : il est d'ordre $8$ et d'indice $3$. Que
ce caractère soit celui par lequel $\mathrm{Aut}(E_a)$ agit sur l'espace des
déformations est notre lecture de « les théorèmes de la rigidification en
car.\ 2 », dont la suite est illisible ; c'est vrai : $\mathrm{Aut}(E_a)$ agit
sur $\omega$ par $u \in \mu_3$, de noyau $Q_8$.} $G'$ agit donc sur les quatre
tangentes possibles par homothéties de $\mathbf{F}_4$, avec deux orbites, $\{0\}$
et $\mathbf{F}_4^{*}$. Comme $G'$ permute transitivement les quatre branches,
l'ensemble de leurs tangentes est une seule orbite, de cardinal divisant $4$ ;
ce ne peut être que $\{0\}$. \emph{Les quatre branches ont en $a$ la même
tangente, la direction horizontale.}\footnote{La conclusion est celle par
laquelle la page 99 s'ouvre (« \ldots{} aux composantes irréductibles sont
égales »), la page 97 s'interrompant sur « Donc les quatre\ldots ». L'argument
d'orbites est le nôtre, construit sur les éléments lisibles des deux pages :
la décomposition $V_1 \times V_2$, l'action via $\mathbf{F}_4^{*}$, la
transitivité. La page ajoute qu'il faudrait encore déterminer lequel des deux
isomorphismes $G'/P' \simeq \mathbf{F}_4^{*}$ intervient ; elle ne le fait
pas.} C'est cohérent avec la page 95 : quatre branches transverses à la fibre,
chacune de multiplicité $1$, donnent la multiplicité $4$ de $M'$ en $p = 2$,
et quatre branches par un même point sont la non-régularité qu'elle y
constatait.

Un dernier paragraphe de la page 99, très raturé, raisonne par l'absurde sur
l'action de $\mu_3$ sur la courbe universelle au-dessus de $M'_3$ : elle
forcerait $\mu_3$ à être central dans $G' \simeq \mathrm{Aut}(E_a)$, « or le
$3$-sous-groupe de Sylow de $G'$ n'est pas central ». Le fait invoqué est
exact (le centre de $\mathrm{SL}(2,\mathbf{F}_3)$ est $\{\pm 1\}$) ; l'hypothèse
réfutée ne se lit pas, et nous ne la restituons pas.\footnote{Une dizaine de
lignes de ce paragraphe, où figurent $\mu_6 = \mu_2\times\mu_3$ et
« $\mathbf{Z}/6\mathbf{Z} \simeq B$ opère », sont biffées et illisibles.}

\end{document}
